\section{Moving sites with fixed masses}
\label{sec:finite-stability}

We now use the moment estimate to control transport potentials along an
optimal coupling of two finite targets. Intermediate sites need only
remain in the enclosing ball \(B(0,L)\); they need not belong to \(Y\).

\begin{proposition}\label{finite:potentials}
Let \(\mu,\nu\) be finitely supported probabilities in \(B(0,L)\).
There are convex \(L\)-Lipschitz finite maxima \(u_0,u_1\) of affine
functions such that
\[
 \nabla u_0=T_\mu,\qquad \nabla u_1=T_\nu
 \quad \rho\text{-almost everywhere},
\]
and
\begin{equation}
 \|u_0-u_1\|_{L^2(\rho)}
 \le (1+\sqrt{C_0})R\,W_2(\mu,\nu),
 \qquad C_0=162.
 \label{eq:finite-potential-bound}
\end{equation}
\end{proposition}

\begin{proof}
Choose an optimal finite coupling of the targets. After deleting zero
weights and combining identical pairs, write it as
\[
 \pi=\sum_{i=1}^N m_i\delta_{(y_i^0,y_i^1)},\qquad
 m_i>0,\quad \sum_i m_i=1.
\]
Finite-dimensional compactness gives existence of such a coupling.
Set
\[
 v_i=y_i^1-y_i^0,\qquad
 y_i(t)=(1-t)y_i^0+ty_i^1,\qquad
 w=W_2(\mu,\nu).
\]
Then
\begin{equation}
 \sum_i m_i|v_i|^2=w^2.
 \label{eq:finite-coupling-speed}
\end{equation}

Optimality implies
\begin{equation}
 (y_i^0-y_j^0)\cdot(y_i^1-y_j^1)\ge0.
 \label{eq:finite-monotonicity}
\end{equation}
Indeed, a negative scalar product would permit swapping equal positive
portions of the masses at these two coupling pairs, decreasing their
quadratic cost. For \(a=y_i^0-y_j^0\) and \(b=y_i^1-y_j^1\),
\[
 |y_i(t)-y_j(t)|^2
 =(1-t)^2|a|^2+t^2|b|^2+2t(1-t)a\cdot b>0
 \qquad(0<t<1).
\]
The last inequality holds because the endpoint pairs are distinct.
Thus the sites are distinct at every interior time, even if some
endpoint sites coincide.

By \Cref{cells:weights}, for \(0<t<1\) there are unique intercepts
\(c_i(t)\), normalized by
\[
 \sum_i m_ic_i(t)=0,
\]
whose cells have masses \(M_i(t)=m_i\). These intercepts are \(C^1\) in
\(t\). Let \(r_i=c_i'(t)\). In particular,
\begin{equation}
 \sum_i m_ir_i=0.
 \label{eq:finite-zero-mean}
\end{equation}
All derivatives in the next calculation are taken at this fixed
interior configuration.

Differentiate the fixed-mass constraints using \Cref{cells:calculus}.
For every intercept direction \(q\), symmetry of the mass and moment
face coefficients gives
\begin{equation}
 \sum_i r_i\,DM_i[q,0]
 =-\sum_i v_i\cdot DB_i[q,0].
 \label{eq:finite-duality}
\end{equation}
To see the symmetry directly, the mass part of
\(\sum_iq_iDM_i[r,v]\) is
\(\sum_{i<j}g_{ij}(q_i-q_j)(r_i-r_j)\). Its site part is
\(\sum_{i<j}(q_i-q_j)b_{ij}\cdot(v_i-v_j)\), which equals
\(\sum_iv_i\cdot DB_i[q,0]\).

The range statement in \Cref{cells:calculus}, together with
\eqref{eq:finite-zero-mean}, permits a choice of \(q\) satisfying
\[
 DM_i[q,0]=m_ir_i\qquad(1\le i\le N).
\]
No bound on this \(q\) is needed. Put \(E_r=\sum_i m_ir_i^2\).
Using \eqref{eq:finite-duality}, weighted Cauchy--Schwarz,
\eqref{eq:finite-coupling-speed}, and \Cref{moments:derivative}, we obtain
\[
 \begin{split}
 E_r
 &\le
 \left(\sum_i m_i|v_i|^2\right)^{1/2}
 \left(\sum_i\frac{|DB_i[q,0]|^2}{m_i}\right)^{1/2}\\
 &\le \sqrt{C_0}Rw
       \left(\sum_i\frac{|m_ir_i|^2}{m_i}\right)^{1/2}
 =\sqrt{C_0}Rw\sqrt{E_r}.
 \end{split}
\]
Consequently
\begin{equation}
 \left(\sum_i m_i|r_i|^2\right)^{1/2}
 \le\sqrt{C_0}Rw.
 \label{eq:finite-intercept-speed}
\end{equation}
This includes \(E_r=0\). If \(N=1\), the normalization gives \(r_1=0\),
and the same conclusion is immediate.

Define
\[
 u_t(x)=\max_i\{x\cdot y_i(t)+c_i(t)\}.
\]
On compact time intervals inside \((0,1)\), this function is uniformly
Lipschitz in time for \(x\in K\). At a strict maximum its time derivative
is \(x\cdot v_i+r_i\). For each fixed time, the exceptional tie set has
zero source measure. Fubini therefore identifies this derivative for
almost every \((t,x)\). Since its active cells have masses \(m_i\),
\[
 \|\partial_tu_t\|_{L^2(\rho)}
 \le R\left(\sum_i m_i|v_i|^2\right)^{1/2}
     +\left(\sum_i m_i|r_i|^2\right)^{1/2}
 \le(1+\sqrt{C_0})Rw.
\]
The fundamental theorem of calculus and Minkowski's inequality imply
\begin{equation}
 \|u_t-u_{t'}\|_{L^2(\rho)}
 \le(1+\sqrt{C_0})Rw\,|t-t'|
 \qquad(0<t,t'<1).
 \label{eq:finite-time-bound}
\end{equation}

It remains to justify the endpoint maps without imposing separation of
their sites. For each fixed \(i\), \eqref{eq:finite-intercept-speed}
gives
\[
 |c_i'(t)|\le\frac{\sqrt{C_0}Rw}{\sqrt{m_i}}.
\]
Hence \(c_i(t)\) has finite limits at \(0\) and \(1\). This individual
estimate is used only for endpoint existence; its mass dependence
does not enter \eqref{eq:finite-time-bound}. The finite maxima converge
uniformly on \(K\) to their endpoint maxima \(u_0,u_1\), so that
inequality extends to the endpoints.

At either endpoint, merge repeated slopes by retaining their largest
intercept. Outside a finite union of hyperplanes there is a unique
maximizing slope. For any such \(x\), all maximizing slopes at
approaching interior times converge to this endpoint slope. Otherwise
a subsequence with a fixed maximizing index would contradict the
strict gap at the endpoint. The interior slope law is
\(\sum_i m_i\delta_{y_i(t)}\), and all slopes are bounded by \(L\).
Bounded convergence against continuous tests shows that the endpoint
slope laws are \(\mu\) and \(\nu\). By \Cref{cells:optimality}, these
slopes are the unique optimal maps to their laws. Taking \(t=0,t'=1\)
in \eqref{eq:finite-time-bound} proves the proposition.
\end{proof}

\begin{remark}[The potential estimate on the cube example]
Return to $K=Y=[-1,1]^d$ and the potentials $u_a,\widetilde u_a$ in
\eqref{sharpness:potentials}, where $0<a<1/2$ and $b=a/2$.
Both potentials have weighted intercept mean $a^2$, so subtracting
$a^2$ realizes the normalization in \Cref{finite:potentials} without
changing their difference. The intermediate maximum
$\max\{-x_1,x_1,a+sbx_2\}$ has the same three cell masses for every
$s\in[0,1]$, so the same constant normalizes the entire path.

For an exact calculation, put $f=\widetilde u_a-u_a$ and
$q=b|t|$. Integrating its square over the $x_1$ slice gives
\[
 \E[f^2\mid x_2=t]=
 \begin{cases}
   aq^2+q^3/3,&t\ge0,\\
   aq^2-2q^3/3,&t<0.
 \end{cases}
\]
For $t\ge0$, this is the contribution $aq^2$ of the old central
interval plus the integral of the squared linear ramp on the added
portion. For $t<0$, the overlap contributes $(a-q)q^2$ and the
removed portion contributes $q^3/3$. Averaging over $t$ yields
\[
 \|u_a-\widetilde u_a\|_{L^2(\rho)}^2
 =\frac{ab^2}{3}-\frac{b^3}{24}
 =\frac{5a^3}{64}
 =\frac5{16}W_2(\mu_a,\nu_a)^2.
\]
Thus the normalized potential distance is proportional to $W_2$,
whereas \eqref{sharpness:map-distance} has map distance of order
$W_2^{1/3}$. The next section recovers this exponent from a general
interpolation inequality for convex gradients.
\end{remark}
